\documentclass[t,aspectratio=169]{beamer}
\usepackage[utf8]{vietnam}
\usepackage{amsthm,amsmath,amssymb}
\usepackage{lmodern}

% Giữ nguyên bố cục và nhận diện của template HUST.
\def\slideContentLeftMargin{0.6cm}
\def\slideContentRightMargin{0.9cm}
\def\hustHeaderHeight{1cm}
\def\slideContentBotMargin{3.5cm}

\usepackage[blue,169]{beamerthemeHUST}
\usetikzlibrary{arrows.meta,positioning}

\renewcommand{\titlePageTitleX}{-0.5cm}
\renewcommand{\titlePageTitleY}{-0.35cm}
\renewcommand{\hustTitleAlign}{\alignLeft}
\renewcommand{\hustSubtitleAlign}{\alignLeft}
\renewcommand{\hustAuthorAlign}{\alignLeft}
\renewcommand{\hustInstAlign}{\alignLeft}

\title{Báo cáo tiến độ tuần 1}
\subtitle{Phân cụm mờ C-Means bán giám sát}
\author{Nguyễn Trí Hiếu}
\institute{Trường Công nghệ Thông tin và Truyền thông (SoICT)\\
Đại học Bách khoa Hà Nội\\
Ngày 29 tháng 9 năm 2026}

\begin{document}

\husttitlepage
\resetPageCount
\showPageNumber
\useStyleHustFull
\useThinBar

% -----------------------------------------------------------------------------
\section{Tiến độ}

\begin{frame}{Mục tiêu và tiến độ tuần 1}
  \begin{columns}[T,onlytextwidth]
    \begin{column}{0.60\textwidth}
      \begin{itemize}
        \item Đã đọc ba bài báo về sSFCM, sSMC-FCM và đánh giá phân cụm.
        \item Đã hiểu tiến trình từ FCM đến hai cơ chế bán giám sát.
        \item Đã đối chiếu ví dụ 20 điểm và các giá trị độ thuộc được báo cáo.
        \item Đã xác định thứ tự hiện thực và thí nghiệm ban đầu.
      \end{itemize}
    \end{column}
    \begin{column}{0.36\textwidth}
      \begin{block}{Kết quả hiện tại}
        \centering
        \textbf{Thuật toán sẽ hiện thực}\\[0.3em]
        \Large\textcolor{hustprimary}{sSMC-FCM}\\[0.5em]
        \normalsize cùng hai baseline FCM và sSFCM
      \end{block}
    \end{column}
  \end{columns}

  \vspace{0.6em}
  \begin{alertblock}{Phạm vi tuần 1}
    Công việc hiện tại mới gồm phân tích bài báo và thiết kế thí nghiệm. Chưa viết mã thuật toán.
  \end{alertblock}
\end{frame}

% -----------------------------------------------------------------------------
\section{Phương pháp}

\begin{frame}{FCM và vai trò của thông tin giám sát}
  \begin{columns}[T,onlytextwidth]
    \begin{column}{0.54\textwidth}
      \begin{block}{FCM chuẩn}
        \[
          J_m(U,V)=\sum_{i=1}^{N}\sum_{k=1}^{C}u_{ik}^{m}D_{ik}^{2},
          \qquad \sum_{k=1}^{C}u_{ik}=1
        \]
        \[
          V_k=\frac{\sum_{i=1}^{N}u_{ik}^{m}x_i}
          {\sum_{i=1}^{N}u_{ik}^{m}}
        \]
      \end{block}
    \end{column}
    \begin{column}{0.42\textwidth}
      \begin{itemize}
        \item $u_{ik}$: độ thuộc của điểm $x_i$ vào cụm $k$.
        \item $V_k$: tâm của cụm $k$.
        \item $m>1$: hệ số mờ hóa dùng chung.
        \item $D_{ik}=\lVert x_i-V_k\rVert$.
      \end{itemize}
    \end{column}
  \end{columns}

  \vspace{0.3em}
  \begin{alertblock}{Vì sao cần bán giám sát?}
    Một số gợi ý điểm-cụm đáng tin cậy có thể định hướng các mẫu khó, trong khi phần dữ liệu còn lại không cần nhãn.
  \end{alertblock}
\end{frame}

\begin{frame}{sSFCM 2009: Độ thuộc giám sát}
  \begin{columns}[T,onlytextwidth]
    \begin{column}{0.52\textwidth}
      \begin{block}{Công thức chính}
        \[
          J(U,V)=\sum_{i=1}^{N}\sum_{k=1}^{C}
          \left|u_{ik}-\bar{u}_{ik}\right|^{m}D_{ik}^{2}
        \]
        \[
          u_{ik}=\bar{u}_{ik}+
          \left(1-\sum_{j=1}^{C}\bar{u}_{ij}\right)
          \frac{D_{ik}^{-2/(m-1)}}
          {\sum_{\ell=1}^{C}D_{i\ell}^{-2/(m-1)}}
        \]
      \end{block}
    \end{column}
    \begin{column}{0.44\textwidth}
      \begin{itemize}
        \item $\bar U$ lưu các độ thuộc được cung cấp trước.
        \item Cặp không có giám sát dùng $\bar u_{ik}=0$.
        \item $\bar U$ tác động trực tiếp đến cập nhật độ thuộc.
        \item Tâm cụm dùng trọng số $|u_{ik}-\bar u_{ik}|^m$.
        \item Thuật toán luân phiên cập nhật độ thuộc và tâm.
      \end{itemize}
    \end{column}
  \end{columns}
\end{frame}

\begin{frame}{sSMC-FCM 2021: Nhiều hệ số mờ hóa}
  \small
  \begin{columns}[T,onlytextwidth]
    \begin{column}{0.47\textwidth}
      \begin{block}{Hàm mục tiêu}
        \[
          J(U,V)=\sum_{i=1}^{N}\sum_{k=1}^{C}
          u_{ik}^{m_{ik}}D_{ik}^{2}
        \]
      \end{block}

      \begin{itemize}
        \item Không cần ma trận độ thuộc giám sát phụ trợ.
        \item Mẫu không giám sát dùng cập nhật FCM chuẩn.
        \item Mẫu có giám sát phải giải Eq. (19), sau đó chuẩn hóa.
      \end{itemize}
    \end{column}
    \begin{column}{0.49\textwidth}
      \begin{block}{Cách giám sát thay đổi $m_{ik}$}
        \centering
        \begin{tikzpicture}[
          node distance=0.35cm,
          box/.style={draw=hustprimary,rounded corners=2pt,fill=hustprimary!7,
            align=center,text width=5.5cm,minimum height=0.75cm,inner sep=4pt},
          arr/.style={-{Latex[length=2mm]},thick,draw=hustprimary}
        ]
          \node[box] (u) {Mẫu không giám sát $x_i$\\$m_{ij}=M$ với mọi cụm $j$};
          \node[box,below=of u] (s) {Mẫu $x_i$ được giám sát vào cụm $k$\\$m_{ik}=M_0>M$, và $m_{ij}=M$ với $j\neq k$};
          \draw[arr] (u) -- node[right,font=\scriptsize]{thêm thông tin cụm đích} (s);
        \end{tikzpicture}
      \end{block}

      \vspace{0.25em}
      \textbf{Tác động:} Hệ số tại cặp được chọn làm thay đổi cả độ thuộc và trọng số cập nhật tâm. Eq. (23) hướng dẫn chọn $M_0$ để đạt $u_{ik}\geq\alpha$. Paper ký hiệu hệ số này là $M'$.
    \end{column}
  \end{columns}
\end{frame}

\begin{frame}{So sánh sSFCM 2009 và sSMC-FCM 2021}
  \centering
  \footnotesize
  \setlength{\tabcolsep}{3pt}
  \renewcommand{\arraystretch}{1.22}
  \begin{tabular}{p{2.15cm}|p{5.55cm}|p{5.55cm}}
    \textbf{Khía cạnh} & \textbf{sSFCM 2009} & \textbf{sSMC-FCM 2021} \\
    \hline
    Giám sát & Các giá trị trong $\bar U$ & Các cặp mẫu-cụm đích \\
    Hàm mục tiêu & $|u_{ik}-\bar u_{ik}|^mD_{ik}^2$ & $u_{ik}^{m_{ik}}D_{ik}^2$ \\
    Tham số chính & Độ thuộc $\bar u_{ik}$ & Hai hệ số $M$ và $M_0>M$ \\
    Cập nhật $U$ & Cộng độ thuộc giám sát rồi phân bổ phần còn lại & FCM cho mẫu không nhãn; giải scalar cho cặp đích \\
    Ảnh hưởng đến tâm & Trọng số $|u_{ik}-\bar u_{ik}|^m$ & Trọng số $u_{ik}^{m_{ik}}$ \\
    Đầu vào thêm & Ma trận phụ trợ $\bar U$ & Cặp mẫu-cụm và $M_0$ \\
    Ý tưởng chính & Thay đổi thông tin độ thuộc & Thay đổi hành vi mờ hóa \\
  \end{tabular}

  \vspace{0.5em}
  \begin{alertblock}{Hệ quả khi lập trình}
    Phương pháp 2021 thay hướng dẫn trực tiếp bằng số mũ theo từng cặp và một bài toán số cho độ thuộc có giám sát.
  \end{alertblock}
\end{frame}

% -----------------------------------------------------------------------------
\section{Kết quả đối chiếu và đánh giá}

\begin{frame}{Ví dụ số 20 điểm dùng chung}
  \centering
  \small
  \setlength{\tabcolsep}{3pt}
  \renewcommand{\arraystretch}{1.22}
  \begin{tabular}{p{4.4cm}|p{3.1cm}|c|c|p{2.0cm}}
    \textbf{Phương pháp} & \textbf{Giám sát} & $\boldsymbol{U_{C1}}$ & $\boldsymbol{U_{C2}}$ & \textbf{Cụm cứng} \\
    \hline
    FCM & không & 0.19 & 0.81 & C2 \\
    sSFCM 2009 & $\bar u_{i1}=0.3$ & 0.45 & 0.55 & C2 \\
    sSFCM 2009 & $\bar u_{i1}=0.6$ & \textbf{0.69} & 0.31 & C1 \\
    sSMC-FCM 2021 & $M=2,\ M_0=4$ & 0.43 & 0.57 & C2 \\
    sSMC-FCM 2021 & $M=2,\ M_0=8$ & \textbf{0.60} & 0.40 & C1 \\
  \end{tabular}

  \vspace{0.6em}
  \begin{block}{Điểm 9 và 10 có cùng độ thuộc được báo cáo}
    Cả hai cơ chế đều kéo dần các mẫu được giám sát về cụm 1. Khi giám sát đủ mạnh, nhãn cứng chuyển từ cụm 2 sang cụm 1.
  \end{block}
  \begin{alertblock}{Giới hạn diễn giải}
    Ví dụ này chỉ chứng minh hành vi tương tự. Nó không đủ để kết luận phương pháp 2021 tốt hơn một cách tổng quát.
  \end{alertblock}
\end{frame}

\begin{frame}{Bộ chỉ số đánh giá ban đầu}
  \small
  \textbf{\textcolor{red}{Quy tắc:}} Paper 2010 đánh giá phân hoạch cứng, vì vậy dùng
  $\hat y_i=\operatorname*{arg\,max}_{k}u_{ik}$ trước khi tính chỉ số.

  \vspace{0.35em}
  \centering
  \footnotesize
  \setlength{\tabcolsep}{3pt}
  \renewcommand{\arraystretch}{1.18}
  \begin{tabular}{p{2.25cm}|p{2.35cm}|p{6.2cm}|p{1.75cm}}
    \textbf{Nhãn thật} & \textbf{Chỉ số} & \textbf{Đo lường} & \textbf{Tốt hơn} \\
    \hline
    Có & ARI & Mức đồng thuận theo cặp, đã hiệu chỉnh ngẫu nhiên & Cao hơn \\
    Có & Jaccard & Đồng thuận trên các cặp được xếp cùng cụm & Cao hơn \\
    Không & SWC / SSWC & Độ chặt và độ tách; SSWC dùng tâm cụm & Cao hơn \\
    Không & VRC & Phân tán giữa cụm so với phân tán trong cụm & Cao hơn \\
  \end{tabular}

  \vspace{0.45em}
  \raggedright
  \small
  \textbf{Lý do chọn:} ARI và Jaccard phù hợp benchmark có nhãn. Silhouette thuộc nhóm mạnh trong paper, còn VRC là phép kiểm tra bổ sung có chi phí thấp.
\end{frame}

% -----------------------------------------------------------------------------
\section{Kế hoạch hiện thực}

\begin{frame}{Pipeline hiện thực và thí nghiệm}
  \begin{columns}[T,onlytextwidth]
    \begin{column}{0.58\textwidth}
      \centering
      \begin{tikzpicture}[
        flow/.style={draw=hustprimary,rounded corners=2pt,fill=hustprimary!7,
          align=center,minimum height=0.58cm,minimum width=2.05cm,
          font=\scriptsize,inner xsep=4pt},
        arr/.style={-{Latex[length=2mm]},thick,draw=hustprimary}
      ]
        \node[flow] (data) {Dữ liệu};
        \node[flow,below=0.34cm of data] (prep) {Tiền xử lý và chia tập giám sát};
        \node[flow,below=0.34cm of prep] (init) {Khởi tạo chung};
        \node[flow,below=0.62cm of init,xshift=-2.35cm] (fcm) {FCM};
        \node[flow,below=0.62cm of init] (ssfcm) {sSFCM};
        \node[flow,below=0.62cm of init,xshift=2.35cm] (ssmc) {sSMC-FCM};
        \node[flow,below=1.08cm of ssfcm] (metric) {Đánh giá và kiểm tra};
        \node[flow,below=0.34cm of metric] (out) {Bảng so sánh và biểu đồ};
        \draw[arr] (data) -- (prep);
        \draw[arr] (prep) -- (init);
        \draw[arr] (init) -- (fcm);
        \draw[arr] (init) -- (ssfcm);
        \draw[arr] (init) -- (ssmc);
        \draw[arr] (fcm) -- (metric);
        \draw[arr] (ssfcm) -- (metric);
        \draw[arr] (ssmc) -- (metric);
        \draw[arr] (metric) -- (out);
      \end{tikzpicture}
    \end{column}
    \begin{column}{0.38\textwidth}
      \begin{block}{Các module dự kiến}
        \small\ttfamily
        src/fcm.py\\
        src/ssfcm.py\\
        src/ssmc\_fcm.py\\
        src/metrics.py\\
        src/datasets.py\\
        src/experiments.py
      \end{block}
      \begin{block}{Thứ tự hiện thực}
        \scriptsize
        \begin{enumerate}
          \item FCM và unit test
          \item Dữ liệu 20 điểm năm 2009
          \item Tái lập sSFCM
          \item sSMC-FCM Eqs. (17)--(20)
          \item Tái lập Table 3 năm 2021
          \item Metrics và dữ liệu lớn hơn
        \end{enumerate}
      \end{block}
    \end{column}
  \end{columns}
\end{frame}

\begin{frame}{Bước tiếp theo và quyết định cần chốt}
  \small
  \begin{columns}[T,onlytextwidth]
    \begin{column}{0.49\textwidth}
      \begin{block}{Các bước coding tiếp theo}
        \begin{enumerate}
          \item Hiện thực FCM và kiểm tra shape, simplex, finite.
          \item Tạo lại bộ 20 điểm và độ thuộc FCM.
          \item Hiện thực sSFCM, tái lập kết quả năm 2009.
          \item Hiện thực solver độ thuộc cho sSMC-FCM.
          \item Tái lập kết quả năm 2021 trước thí nghiệm lớn.
        \end{enumerate}
      \end{block}
    \end{column}
    \begin{column}{0.47\textwidth}
      \begin{alertblock}{Các câu hỏi cần trao đổi}
        \begin{itemize}
          \item Tolerance, bước tăng và cận trên cho Eq. (19).
          \item Khởi tạo tâm và norm dừng.
          \item Xử lý khoảng cách 0 và cụm rỗng.
          \item Dataset thật và tỷ lệ mẫu giám sát.
        \end{itemize}
      \end{alertblock}
      \textbf{Cổng tuần 2:} tái lập độ thuộc của điểm 9/10 trước khi thử dữ liệu mới.
    \end{column}
  \end{columns}
\end{frame}

% -----------------------------------------------------------------------------
\section{Tài liệu tham khảo}
\useStyleLogoOnly

\begin{frame}{Tài liệu tham khảo}
  \small
  \begin{thebibliography}{9}
    \bibitem{endo2009}
    Y. Endo, Y. Hamasuna, M. Yamashiro, and S. Miyamoto,
    ``On Semi-Supervised Fuzzy c-Means Clustering,''
    \emph{IEEE International Conference on Fuzzy Systems}, 2009.

    \bibitem{khang2021}
    T. D. Khang, M.-K. Tran, and M. Fowler,
    ``A Novel Semi-Supervised Fuzzy C-Means Clustering Algorithm Using Multiple Fuzzification Coefficients,''
    \emph{Algorithms}, vol. 14, no. 9, article 258, 2021.
    DOI: 10.3390/a14090258.

    \bibitem{vendramin2010}
    L. Vendramin, R. J. G. B. Campello, and E. R. Hruschka,
    ``Relative Clustering Validity Criteria: A Comparative Overview,''
    \emph{Statistical Analysis and Data Mining}, vol. 3, pp. 209--235, 2010.
    DOI: 10.1002/sam.10080.
  \end{thebibliography}
\end{frame}

\hidePageNumber
\hustcontactpage{\hfill CÂU HỎI VÀ TRAO ĐỔI \hfill}{
  \textbf{Mục tiêu tuần 2:} tái lập đúng các giá trị độ thuộc đã công bố trước khi mở rộng thí nghiệm.
}

\end{document}
